% !TeX root = ../tesi.tex
% !TeX encoding = UTF-8
% !TeX spellcheck = it_IT

\chapter{Background e lavori correlati}
\label{chap:background}

In questo capitolo presento TruFor, il sistema che utilizzo per studiare le manipolazioni. TruFor è un framework per l'analisi \emph{forense} delle immagini, cioè l'esame delle tracce lasciate dall'acquisizione e dalle elaborazioni successive, per individuare possibili manipolazioni. Introduco i compiti dell'analisi forense e descrivo il funzionamento di TruFor. Esamino poi i risultati e i limiti che gli autori riportano nel loro lavoro. Uso questi elementi per motivare la scelta del sistema nello studio della comunicazione politica sui social media.

Presento poi gli strumenti statistici per l'analisi delle associazioni e i descrittori visivi che adopero per rappresentare le caratteristiche dell'immagine. Infine introduco gli strumenti per il riconoscimento ottico dei caratteri, e per descrivere il contenuto delle regioni sospette.

\section{TruFor}
\label{sec:background-trufor}

\subsection{Rilevamento e localizzazione delle manipolazioni}
\label{sec:background-rilevamento-localizzazione}

Isolo due compiti dell'analisi. Il \emph{rilevamento} classifica un'immagine come autentica o manipolata. La \emph{localizzazione} indica le regioni che il sistema segnala come manipolate. Descrivo questo compito con una \emph{segmentazione binaria}: il sistema assegna ciascun pixel alla classe ``autentico'' o ``manipolato''~\cite{Guillaro_2023_CVPR}.

Osservo manipolazioni come il copia-incolla e l'inserimento di contenuti sintetici.

Valuto positiva la classe ``manipolato'' e negativa la classe ``autentico''. Questa convenzione vale sia per le immagini nel rilevamento sia per i pixel nella localizzazione. TruFor può segnalare come manipolate anche regioni autentiche, producendo falsi positivi a livello di pixel. Nel rilevamento, i falsi positivi possono costituire una parte consistente delle segnalazioni quando le manipolazioni sono poco frequenti. Quando interpreto i risultati, distinguo quindi gli errori di localizzazione dagli errori di rilevamento.

\subsection{Indizi visivi e impronta Noiseprint++}
\label{sec:background-noiseprint}

L'immagine RGB descrive il contenuto visibile; l'\emph{estrattore Noiseprint++} genera un'impronta, \emph{Noiseprint++}, che evidenzia le incoerenze tra le tracce presenti nelle diverse regioni. Trufor combina la fotografia con la sua impronta.

Noiseprint++ estende \emph{Noiseprint}, un'impronta che rappresenta le tracce dei processi interni alla fotocamera. Noiseprint++ include anche le tracce delle elaborazioni successive all'acquisizione. L'addestramento dell'estrattore di impronta avviene con immagini autentiche. L'addestramento è auto-supervisionato: gli autori ricavano il segnale di apprendimento dalle relazioni tra gli esempi.

Gli autori utilizzano 24.700 immagini acquisite con 1.475 modelli di fotocamera. Ne estraggono porzioni di $64 \times 64$ pixel. L'addestramento segue un criterio \emph{contrastivo}: produce impronte simili nelle aree che condividono le stesse proprietà e diverse nelle aree che differiscono. Le proprietà richieste sono il modello di fotocamera, la posizione nello spazio e la storia delle elaborazioni. Gli sviluppatori simulano quest'ultima con 512 combinazioni di ridimensionamento, compressione e variazioni di illuminazione.

\subsection{Architettura e apprendimento}
\label{sec:background-trufor-architettura}

Il sistema elabora l'RGB e Noiseprint++ attraverso due rami. Ogni ramo estrae delle \emph{caratteristiche}, rappresentazioni numeriche delle informazioni della figura. La fusione delle caratteristiche usa una rete di tipo \emph{Transformer}, che apprende relazioni tra gli elementi delle rappresentazioni.

L'\emph{encoder} estrae le caratteristiche che alimentano due \emph{decoder}. Ciascun decoder produce una mappa con la stessa risoluzione della fotografia.

Il decoder delle anomalie produce la \emph{mappa delle anomalie}, che assegna a ciascun pixel un valore che esprime quanto il sistema lo considera sospetto. I valori alti indicano le regioni segnalate come manipolate. Pixel autentici possono ricevere valori alti e generare falsi positivi.

Il decoder della confidenza produce la \emph{mappa di confidenza}, che stima l'affidabilità della previsione della mappa delle anomalie per ciascun pixel. Una confidenza alta indica che il sistema considera affidabile la previsione. La confidenza riguarda entrambe le classi: ``autentico'' e ``manipolato''.

Riassumo nella Figura~\ref{fig:trufor-mappe} il percorso dall'immagine alle due mappe.

\begin{figure}[htbp]
    \centering
    \begin{tikzpicture}[
        blocco/.style={draw=black!65, rounded corners=2pt, align=center,
            font=\small, text width=2.5cm, minimum height=0.85cm, inner sep=5pt},
        dato/.style={blocco, fill=black!4},
        modulo/.style={blocco, fill=sapred!20},
        flusso/.style={-{Stealth[length=2mm]}, semithick, draw=black!70}
    ]
        \node[dato] (rgb) at (-3.5,0) {Immagine RGB};
        \node[modulo] (estrattore) at (0,0) {Estrattore\\Noiseprint++};
        \node[dato] (impronta) at (3.5,0) {Impronta\\Noiseprint++};
        \node[modulo, text width=4.4cm] (encoder) at (0,-1.6)
            {Encoder a due rami\\e fusione delle caratteristiche};
        \node[modulo] (anomalie) at (-2,-3.1) {Decoder\\delle anomalie};
        \node[modulo] (confidenza) at (2,-3.1) {Decoder\\della confidenza};
        \node[dato] (mappaa) at (-2,-4.4) {Mappa\\delle anomalie};
        \node[dato] (mappac) at (2,-4.4) {Mappa\\di confidenza};
        \draw[flusso] (rgb) -- (estrattore);
        \draw[flusso] (estrattore) -- (impronta);
        \draw[flusso] (rgb.south) |- (encoder.west);
        \draw[flusso] (impronta.south) |- (encoder.east);
        \draw[flusso] (encoder.south) -- ++(0,-0.15) -| (anomalie.north);
        \draw[flusso] (encoder.south) -- ++(0,-0.15) -| (confidenza.north);
        \draw[flusso] (anomalie) -- (mappaa);
        \draw[flusso] (confidenza) -- (mappac);
    \end{tikzpicture}
    \caption{Dall'immagine alle mappe di TruFor. I riquadri colorati rappresentano i componenti appresi; quelli grigi rappresentano i dati. Schema semplificato dalla Figura~2~\cite{Guillaro_2023_CVPR}.}
    \label{fig:trufor-mappe}
\end{figure}

\newpage
Gli autori addestrano TruFor in tre fasi, che riassumo nella Tabella~\ref{tab:trufor-addestramento}. La prima usa il criterio contrastivo descritto nella Sezione~\ref{sec:background-noiseprint}. Nella seconda usano un apprendimento \emph{supervisionato}: gli autori accompagnano gli esempi con le risposte di riferimento. Per ogni istantanea una \emph{maschera di riferimento}, o \emph{ground truth}, indica la classe ``autentico'' o ``manipolato'' di ciascun pixel.

Nella terza fase il decoder della confidenza apprende come stimare la \emph{True Class Probability}, la probabilità che il modello assegna alla classe corretta di ciascun pixel. Questa quantità assume valori alti quando la previsione favorisce la classe corretta. Il decoder della confidenza apprende così come stimare l'affidabilità della mappa delle anomalie.

\begin{table}[htbp]
    \centering
    \small
    \caption{Le tre fasi di addestramento di TruFor~\cite{Guillaro_2023_CVPR}.}
    \label{tab:trufor-addestramento}
    \begin{tabular}{@{}p{0.07\linewidth}p{0.35\linewidth}p{0.49\linewidth}@{}}
        \toprule
        Fase & Componenti addestrati & Obiettivo \\
        \midrule
        1 & Estrattore Noiseprint++ & Produrre impronte simili per porzioni con lo stesso modello di fotocamera, posizione e storia delle elaborazioni. \\
        \addlinespace
        2 & Encoder e decoder delle anomalie & Assegnare a ciascun pixel la classe indicata dalla maschera di riferimento. \\
        \addlinespace
        3 & Decoder della confidenza e classificatore & Stimare l'affidabilità delle previsioni per pixel e riconoscere le immagini manipolate. \\
        \bottomrule
    \end{tabular}
\end{table}

\subsection{Output e loro interpretazione}
\label{sec:background-trufor-output}

TruFor restituisce la mappa delle anomalie e la mappa di confidenza, definite nella Sezione~\ref{sec:background-trufor-architettura}, insieme a un punteggio globale.

La visualizzazione della mappa delle anomalie come mappa di calore evidenzia le regioni da esaminare. Interpreto queste regioni insieme alla mappa di confidenza, per distinguere le segnalazioni che il sistema considera affidabili da quelle che richiedono cautela.

Nel confronto con la letteratura, gli scienziati osservano che numerosi metodi ricavano la valutazione dell'intera fotografia dalle previsioni sui singoli pixel. Gli errori nelle previsioni locali possono così portare a segnalare come manipolata un'immagine autentica. TruFor usa la mappa di confidenza per ridurre il peso delle previsioni meno affidabili.

Il \emph{punteggio globale}, o \emph{integrity score}, esprime la valutazione sul media nel suo complesso. TruFor calcola le statistiche, come media e massimo, delle mappe attraverso il \emph{pooling}, un'operazione di aggregazione. Il classificatore usa queste informazioni per calcolare il punteggio globale.

Nell'interpretazione seguo la priorità proposta dagli autori: considero prima il punteggio globale e poi esamino le mappe dei casi segnalati.

\subsection{Risultati riportati e confronto con la letteratura}
\label{sec:background-trufor-risultati}

Esamino i risultati per valutare due proprietà: la \emph{generalizzazione}, la capacità di riconoscere manipolazioni diverse da quelle di addestramento, e la \emph{robustezza}, la capacità di mantenere le prestazioni dopo trasformazioni dell'immagine. Gli esperimenti comprendono fotoritocchi e manipolazioni locali prodotte con modelli generativi.

Confronto TruFor con \emph{CAT-Net v2}, una rete neurale che combina il contenuto RGB e le tracce della compressione JPEG per localizzare le manipolazioni. Negli esperimenti CAT-Net v2 occupa il secondo posto per F1-score medio a livello di pixel. Riporto il confronto nella Tabella~\ref{tab:trufor-medie}~\cite{Guillaro_2023_CVPR}.

Per la localizzazione, l'F1-score combina precisione e richiamo a livello di pixel. Per il rilevamento, l'AUC (Area Under the Curve) misura l'area sotto la curva ROC (Receiver Operating Characteristic), che mette in relazione i tassi di veri e falsi positivi al variare della soglia. L'accuratezza bilanciata misura invece la media delle percentuali di riconoscimento delle due classi. Valori maggiori indicano prestazioni migliori per tutte e tre le metriche.

\begin{table}[htbp]
    \centering
    \small
    \caption{Prestazioni medie riportate: localizzazione su otto dataset e rilevamento su sette. Fonte: Tabelle~1 e~2~\cite{Guillaro_2023_CVPR}.}
    \label{tab:trufor-medie}
    \begin{tabular}{@{}llrr@{}}
        \toprule
        Compito e metrica & Soglia & TruFor & CAT-Net v2 \\
        \midrule
        Localizzazione: F1-score & Ottimale per immagine & $0{,}785$ & $0{,}709$ \\
        Localizzazione: F1-score & Fissa a $0{,}5$ & $0{,}696$ & $0{,}601$ \\
        Rilevamento: AUC & Variabile & $0{,}857$ & $0{,}797$ \\
        Rilevamento: acc. bilanciata & Fissa a $0{,}5$ & $0{,}781$ & $0{,}649$ \\
        \bottomrule
    \end{tabular}
\end{table}

TruFor supera CAT-Net v2 nelle quattro medie. Separo i risultati con soglia fissa da quelli con soglia ottimale: gli autori scelgono quest'ultima massimizzando l'F1-score per ogni figura, usando la maschera di riferimento. La soglia fissa permette invece di valutare le previsioni puntuali senza questa scelta specifica.

Distinguo due dataset: \emph{OpenForensics} contiene manipolazioni locali di volti basate su reti generative avversarie GAN (Generative Adversarial Networks); \emph{CocoGlide} include manipolazioni locali prodotte con \emph{GLIDE}, un modello a diffusione. La Tabella~\ref{tab:trufor-generative} mostra i risultati per questi dataset.

\begin{table}[htbp]
    \centering
    \small
    \caption{F1-score di localizzazione per dataset e metodo. Il grassetto indica il valore maggiore fra i due metodi per ciascun dataset e soglia. Fonte: Tabella~1~\cite{Guillaro_2023_CVPR}.}
    \label{tab:trufor-generative}
    \begin{tabular}{@{}llrr@{}}
        \toprule
        Dataset e generazione & Metodo & \shortstack{Soglia ottimale\\per immagine} & \shortstack{Soglia fissa\\a $0{,}5$} \\
        \midrule
        OpenForensics (GAN) & TruFor & $0{,}901$ & $0{,}827$ \\
                            & CAT-Net v2 & $\mathbf{0{,}947}$ & $\mathbf{0{,}899}$ \\
        \addlinespace
        CocoGlide (diffusione) & TruFor & $\mathbf{0{,}720}$ & $\mathbf{0{,}523}$ \\
                              & CAT-Net v2 & $0{,}603$ & $0{,}434$ \\
        \bottomrule
    \end{tabular}
\end{table}

CAT-Net v2 supera TruFor su OpenForensics; TruFor supera CAT-Net v2 su CocoGlide. Il vantaggio medio di TruFor comprende quindi risultati diversi sui singoli dataset. Questi esperimenti valutano l'individuazione delle regioni manipolate su due raccolte diverse: il confronto diretto tra GAN e diffusione richiede di controllare anche i contenuti di partenza e le manipolazioni.

Gli autori valutano anche immagini transitate su Facebook e WhatsApp. Il confronto include CAT-Net v2, IF-OSN e MVSS-Net. \emph{IF-OSN} è un metodo progettato per individuare manipolazioni sui social. Il suo addestramento simula le degradazioni introdotte dalle piattaforme. \emph{MVSS-Net} è una rete neurale che combina indizi del rumore e dei bordi delle regioni manipolate per il rilevamento e l'individuazione delle aree interessate.

Nei tre dataset della Tabella~4, TruFor supera questi metodi nell'F1-score con soglia fissa~\cite{Guillaro_2023_CVPR}. Questi risultati motivano la scelta di TruFor per le immagini politiche sui social media. Le prestazioni riportate riguardano i dataset dei creatori di TruFor; la valutazione sul dataset della tesi resta un passaggio distinto.

\subsection{Ambito di applicazione e limiti}
\label{sec:background-trufor-limiti}

Gli autori circoscrivono le capacità di TruFor alle manipolazioni locali. Il rilevamento delle immagini interamente generate resta fuori dalle capacità del sistema. La ricerca di incoerenze locali richiede regioni autentiche come riferimento. Differenzio quindi l'inserimento di contenuti sintetici in una fotografia dalla generazione totale.

L'addestramento richiede maschere di riferimento per i pixel e procede in fasi separate. Gli autori propongono come direzione di ricerca un addestramento congiunto per i componenti dell'intero sistema. La proposta include l'uso di etichette che indicano la classe dell'intera immagine.

Il confronto con la letteratura fornisce due motivi per scegliere TruFor. Rispetto ai metodi che ricavano la decisione dalla sola mappa delle anomalie, TruFor considera anche l'affidabilità delle previsioni. Rispetto a CAT-Net v2, ottiene valori medi superiori nelle metriche di localizzazione e rilevamento. Queste evidenze riguardano i dataset del confronto e motivano la scelta di TruFor per la mia analisi. Uso i suoi output per esaminare le regioni sospette delle immagini politiche. Studio poi la relazione tra le caratteristiche visive e l'\emph{engagement}, cioè le interazioni degli utenti con il post.

\newpage
\section{Strumenti statistici per l'analisi delle associazioni}
\label{sec:background-misure-statistiche}

Per studiare le relazioni tra il punteggio di TruFor, le caratteristiche visive e l'engagement impiego correlazioni e modelli statistici.

Considero il coefficiente di Pearson per le associazioni lineari~\cite{10.1098/rspl.1895.0041} e quello di Spearman per le associazioni monotone, basate sull'ordinamento dei valori~\cite{ca468a70-0be4-389a-b0b9-5dd1ff52b33f}. Per interpretare i test statistici considero anche il valore di $p$, seguendo le indicazioni dell'American Statistical Association: la significatività statistica va distinta dall'intensità dell'associazione e dalla sua rilevanza~\cite{Wasserstein02042016}.

Per descrivere l'engagement continuo uso la regressione lineare, stimata con il metodo dei minimi quadrati ordinari, OLS (Ordinary Least Squares). Ne valuto l'adattamento ai dati attraverso il coefficiente di determinazione $R^2$~\cite{James2013}.

Per l'engagement rappresentato da due classi sfrutto invece un GLM (Generalized Linear Model) binomiale con collegamento logit~\cite{68aee965-a8a0-3e72-9f89-8d89ae91a62b}. Lo pseudo-$R^2$ di McFadden permette di confrontarlo con un modello che contiene la sola intercetta~\cite{mcfadden1972conditional}.

Calcolo le correlazioni con SciPy~\cite{2020SciPy-NMeth} e stimo i modelli con statsmodels~\cite{josef_perktold_2024_10984387}. Nella Sezione~\ref{sec:materiali-misure-statistiche} presento le definizioni e le formule delle misure e dei modelli adottati.

\section{Descrittori visivi dell'immagine}
\label{sec:background-feature-visive}

Per studiare la relazione tra l'aspetto delle immagini e l'engagement uso delle caratteristiche visive. Considero tono, esposizione, fuoco, nitidezza, colore, distribuzione delle intensità, composizione e struttura.

I lavori di Ke e colleghi e di Datta e colleghi mostrano come descrittori di questo tipo possano rappresentare proprietà rilevanti per la qualità e l'estetica fotografica~\cite{1640788,10.1007/11744078_23}. Nel mio studio li uso per descrivere le immagini e confrontarle con i punteggi di engagement e di TruFor.

Per rappresentare l'intensità riprendo i pesi dei canali di colore della raccomandazione ITU-R BT.601~\cite{itu_bt601}.

Le misure di dettaglio si collegano agli studi sulla messa a fuoco di Pech-Pacheco e colleghi e alla rassegna degli operatori di misura di Pertuz e colleghi~\cite{903548,PERTUZ20131415}. Per la sfocatura seguo il principio di Crété-Roffet e colleghi: confronto le variazioni tra pixel vicini prima e dopo un ulteriore filtraggio~\cite{creteroffet:hal-00232709}.

Per il colore uso la misura di vivacità di Hasler e Süsstrunk~\cite{10.1117/12.477378}. Il lavoro di Jonauskaite e Thorstenson offre un contesto sul rapporto tra colore ed emozioni~\cite{Jonauskaite2024}; definisco invece il calore della pipeline attraverso il bilanciamento dei canali rosso e blu. Per la distribuzione delle intensità uso anche l'entropia di Shannon~\cite{6773024}.

Gli operatori di Sobel e il rilevamento dei bordi di Canny forniscono gli strumenti per descrivere le variazioni locali e la struttura dell'immagine~\cite{996,4767851}.

Eseguo le elaborazioni delle immagini con OpenCV~\cite{opencv_library} e i calcoli di asimmetria e curtosi con SciPy~\cite{2020SciPy-NMeth}. 

Nella Sezione~\ref{sec:materiali-feature-visive} definisco le 19 caratteristiche e le scelte di calcolo adottate nel progetto, distinguendole dai principi generali dei lavori citati.

\section{Descrizione del contenuto delle regioni sospette}
\label{sec:background-descrizioni}

La localizzazione indica quali pixel esaminare. Per studiare l'associazione tra il loro contenuto e l'engagement, costruisco una rappresentazione testuale delle regioni selezionate. Combino due compiti: l'\emph{image captioning}, che genera una didascalia del contenuto visivo, e il riconoscimento ottico dei caratteri, o OCR (Optical Character Recognition), che trascrive le scritte presenti nell'immagine.

\subsection{BLIP e la generazione delle didascalie}
\label{subsec:background-blip}

BLIP (Bootstrapping Language-Image Pre-training) apprende da coppie immagine-testo. L'encoder visivo è un Vision Transformer: divide l'immagine in patch e le rappresenta come una sequenza di vettori. L'architettura \emph{MED (Multimodal mixture of Encoder-Decoder)} combina questa rappresentazione con tre modalità testuali: codifica del solo testo, codifica del testo insieme all'immagine e generazione di testo condizionata dall'immagine~\cite{pmlr-v162-li22n}.

Per distinguere le ultime due modalità considero il funzionamento dell'attenzione. L'\emph{attenzione incrociata} permette al componente testuale di pesare le informazioni visive in funzione del testo elaborato. L'encoder testuale usa un'attenzione bidirezionale, che considera anche le parole successive. Il decoder usa invece un'attenzione causale: quando produce una nuova unità può consultare soltanto quelle già disponibili. Encoder e decoder condividono parte dei parametri, ma mantengono separati gli strati di attenzione interna.

La generazione procede un \emph{token}, un'unità del vocabolario del modello, alla volta. Un token può corrispondere a una parte di parola. Indico con $J$ l'immagine fornita al modello e con $s_1,\ldots,s_L$ i token di una didascalia, incluso il simbolo finale. La sua probabilità si scompone come:
\begin{equation}
    P(s_1,\ldots,s_L\mid J)
    =\prod_{j=1}^{L}P(s_j\mid s_1,\ldots,s_{j-1},J).
    \label{eq:blip-generazione}
\end{equation}

Il preaddestramento combina i tre obiettivi della Tabella~\ref{tab:blip-obiettivi}. Le modalità di codifica insegnano al modello a confrontare immagini e testi; quella generativa gli insegna a produrre una descrizione.

\begin{table}[htbp]
    \centering
    \small
    \caption{Obiettivi del preaddestramento BLIP e informazioni apprese~\cite{pmlr-v162-li22n}.}
    \label{tab:blip-obiettivi}
    \begin{tabular}{@{}>{\raggedright\arraybackslash}p{0.34\linewidth}>{\raggedright\arraybackslash}p{0.59\linewidth}@{}}
        \toprule
        Obiettivo & Funzione \\
        \midrule
        Image-Text Contrastive (ITC) & Avvicinare i vettori di immagini e testi corrispondenti e separarli da quelli degli altri abbinamenti. \\
        \addlinespace
        Image-Text Matching (ITM) & Classificare una coppia immagine-testo come corrispondente o meno, analizzando insieme le due modalità. \\
        \addlinespace
        Language Modeling (LM) & Predire il token successivo della didascalia usando l'immagine e i token precedenti. \\
        \bottomrule
    \end{tabular}
\end{table}

Gli autori affrontano anche il rumore nelle coppie raccolte dal web con CapFilt (Captioning and Filtering). Adattano un generatore e un filtro su coppie annotate di COCO. Il generatore aggiunge didascalie sintetiche; il filtro scarta i testi giudicati incoerenti, sia originali sia generati. Le coppie conservate, unite a quelle annotate, servono per preaddestrare un nuovo modello.

Nel progetto uso il modello già addestrato per ottenere una descrizione del contenuto selezionato dalla maschera. Nella mia pipeline eseguo soltanto l'inferenza. Una didascalia può non cogliere tutti i dettagli: il testo prodotto va interpretato come una stima del contenuto visivo.

\subsection{OCR: localizzazione e riconoscimento del testo}
\label{subsec:background-ocr}

Una didascalia può descrivere una persona o un manifesto senza riportarne le scritte. Affianco quindi a BLIP \emph{EasyOCR}, una libreria che combina un rilevatore delle regioni testuali e un riconoscitore dei caratteri. La documentazione ufficiale identifica CRAFT e CRNN come riferimenti per questi due componenti~\cite{easyocr_software}.

\emph{CRAFT (Character Region Awareness for Text detection)} è una rete convoluzionale che produce due mappe: la \emph{mappa delle regioni} localizza i caratteri e la \emph{mappa di affinità} evidenzia i collegamenti tra caratteri adiacenti della stessa parola. Applicando soglie alle mappe, il metodo individua le \emph{componenti connesse}, cioè insiemi contigui di pixel, e le racchiude in riquadri. Queste regioni indicano dove si trova il testo, che viene poi trascritto dal riconoscitore~\cite{DBLP:journals/corr/abs-1904-01941}.

La \emph{CRNN (Convolutional Recurrent Neural Network)} trasforma l'area testuale in una stringa attraverso tre passaggi. Gli strati convoluzionali estraggono mappe di caratteristiche, lette per colonne da sinistra a destra per formare una sequenza di vettori. Gli strati ricorrenti \emph{LSTM (Long Short-Term Memory)} bidirezionali elaborano la sequenza nei due versi, usando il contesto a sinistra e a destra. Infine, per ogni elemento, la rete produce una distribuzione sui caratteri possibili e su un simbolo speciale di assenza di carattere, detto \emph{blank}~\cite{DBLP:journals/corr/ShiBY15}.

La trascrizione usa il criterio \emph{CTC (Connectionist Temporal Classification)}. Più sequenze di predizioni possono rappresentare la stessa parola: per ottenere la stringa finale si accorpano prima i simboli uguali consecutivi e poi si eliminano i blank~\cite{10.1145/1143844.1143891}.

Indico il blank con $\varnothing$: questo simbolo permette di separare due lettere uguali che devono essere mantenute. Infatti, $(\mathrm{l},\varnothing,\mathrm{l})$ diventa \texttt{ll}, mentre $(\mathrm{l},\mathrm{l})$ diventa \texttt{l}. Per esempio, la sequenza seguente produce la parola \texttt{voto}:
\[
    (\mathrm{v},\mathrm{v},\varnothing,\mathrm{o},\mathrm{o},\mathrm{t},\mathrm{o})
    \;\longrightarrow\;
    \texttt{voto}.
\]

Durante l'addestramento, il CTC somma le probabilità delle sequenze compatibili con la trascrizione di riferimento. La rete può quindi apprendere da coppie ritaglio-stringa senza conoscere in anticipo la posizione di ogni carattere.

EasyOCR restituisce per ogni regione testuale le coordinate del riquadro che la delimita, una trascrizione e un punteggio di confidenza~\cite{easyocr_software}. Uso quest'ultimo per filtrare le letture. La confidenza riguarda il riconoscimento del testo ed è distinta dalla confidenza forense di TruFor.

BLIP e EasyOCR forniscono informazioni complementari: in un manifesto, BLIP può descrivere il ritratto e lo sfondo, mentre l'OCR legge cognome e slogan. Entrambi elaborano i pixel conservati dalla maschera, dalla cui selezione dipende l'utilità del testo prodotto.
